% GENERATED FILE. Edit canonical lessons, then run npm run generate:books.
% canonical-source-hash: fnv1a64:4d5c7c628cccdf90
% canonical-lessons: AR-C05-naam, AR-C05-la

\chapter{Yes and No}
\label{ch:ar-yes-and-no}
\hlchaptermodality{5}

\begin{chapteropening}
\textbf{By the end of this chapter:} \emph{I can answer a question in Arabic with naʿam or lā, and write lā as the single joined lām-alif shape.}

Say the yes/no pair and form lā as one lām-alif ligature rather than two separate letters.
\end{chapteropening}

\section[\ar{نعم}]{\ar{نعم} (naʿam) — yes}

\label{lesson:AR-C05-naam}



\begin{quote}
A short word to mean yes — with one sound English simply doesn't have.
\end{quote}

\subsection*{You'll want to know — The word}
\textbf{\ar{نعم}} = \textbf{yes}, said \textbf{naʿam}. Arabic reads \textbf{right to left}, so its three letters run:

\begin{quote}
\ar{ن}  (n) · \ar{ع}  (ʿ) · \ar{م}  (m)  $\to$  \textbf{na‑ʿa‑m}
\end{quote}

The middle letter \textbf{\ar{ع}} is the famous \textbf{ʿayn} — a \emph{pharyngeal} sound made deep in the throat, a tight squeeze with no English equivalent. Don't worry about mastering it today; just know that the little apostrophe-like mark in \emph{naʿam} (the \textbf{ʿ}) is standing in for a real, throaty consonant, not a pause.

\subsection*{You'll want to know — Saying yes, and other ways}
\emph{naʿam} is the clear, standard \textbf{yes} — polite and safe anywhere. Two neighbours you'll also hear:

\begin{itemize}
\raggedright
  \item \textbf{\ar{أجل}} (\emph{ajal}) — a more formal, emphatic \textquotedblleft{}yes, indeed.\textquotedblright{}
  \item In everyday speech, many dialects say \textbf{\ar{أيوة}} (\emph{aywa}) for a casual \textquotedblleft{}yeah.\textquotedblright{}
\end{itemize}

And a useful twist: said as a question, \textbf{naʿam?} also means \textbf{\textquotedblleft{}pardon? / you were saying?\textquotedblright{}} — the same stretch English gives \textquotedblleft{}yes?\textquotedblright{}

\subsection*{Your turn}
Say these aloud:

\begin{itemize}
\raggedright
  \item \textquotedblleft{}naʿam\textquotedblright{} — feel the \ar{ع} squeeze in the middle of the throat
  \item right to left — \ar{ن} · \ar{ع} · \ar{م}
  \item as a question — \textquotedblleft{}naʿam?\textquotedblright{} = \textquotedblleft{}pardon?\textquotedblright{}
\end{itemize}

\begin{tcolorbox}[breakable,colback=teal!4,colframe=teal!35!black,title={Before you move on}]
How do you say yes in Arabic? (\textbf{naʿam}, \ar{نعم}.) Which way does the script run? (\textbf{Right to left}.) What is the \textbf{\ar{ع}} in the middle — a pause or a sound? (A \textbf{sound} — the throaty \emph{ʿayn}, with no English match.) What does \emph{naʿam?} mean when you say it as a question? (\textbf{\textquotedblleft{}Pardon?\textquotedblright{}})
\end{tcolorbox}
\section[\ar{لا}]{\ar{لا} (lā) — no}

\label{lesson:AR-C05-la}



\begin{quote}
Two letters, one long vowel — and one of the most recognisable words in the language.
\end{quote}

\subsection*{You'll want to know — The word}
\textbf{\ar{لا}} = \textbf{no}, said \textbf{lā} (a long \emph{aah}). It's just two letters — \textbf{\ar{ل}} (l) + \textbf{\ar{ا}} (the long \emph{ā}, \emph{alif}) — but they're written joined into one sweeping shape, \textbf{\ar{لا}}, a \emph{lam‑alif} that Arabic always writes as a single unit. It's one of the first letter-joins learners recognise on sight.

\subsection*{You'll want to know — More than an answer}
\textbf{lā} isn't only the reply \textquotedblleft{}no.\textquotedblright{} It is Arabic's great word of \textbf{negation}, and it opens the most-spoken sentence in the language — the first half of the \emph{shahada}:

\begin{quote}
\textbf{\ar{لا} \ar{إله} \ar{إلا} \ar{الله}} — \emph{lā ilāha illā llāh} — \textquotedblleft{}\textbf{There is no} god but God.\textquotedblright{}
\end{quote}

That opening \textbf{lā} is your word. The same \emph{lā} also negates present-tense verbs — \textbf{\ar{لا} \ar{أعرف}} \emph{lā aʿrif}, \textquotedblleft{}I do \textbf{not} know\textquotedblright{} — so, like Spanish \emph{no}, one little word is both the answer \textquotedblleft{}no\textquotedblright{} and the sentence's \textquotedblleft{}not.\textquotedblright{}

\begin{grammarlens}[title={Grammar Lens — Yes and no, together}]
You now have the pair:

\noindent
\begin{tabularx}{\linewidth}{@{}>{\raggedright\arraybackslash}X>{\raggedright\arraybackslash}X>{\raggedright\arraybackslash}X@{}}
\toprule
\textbf{} & \textbf{Arabic} & \textbf{said} \\
\midrule
yes & \textbf{\ar{نعم}} & \emph{naʿam} \\
no & \textbf{\ar{لا}} & \emph{lā} \\
\bottomrule
\end{tabularx}
\end{grammarlens}

\subsection*{Your turn}
Say these aloud:

\begin{itemize}
\raggedright
  \item \textquotedblleft{}lā\textquotedblright{} — a long, open \emph{aah}
  \item the join — \ar{ل} + \ar{ا} written as one, \ar{لا}
  \item the pair — \textquotedblleft{}naʿam / lā\textquotedblright{}
\end{itemize}

\begin{tcolorbox}[breakable,colback=teal!4,colframe=teal!35!black,title={Before you move on}]
How do you say no in Arabic? (\textbf{lā}, \ar{لا}.) What two letters make it, and how are they written? (\textbf{\ar{ل}} + \textbf{\ar{ا}}, joined as a single \emph{lam‑alif} \ar{لا}.) Which famous sentence does \emph{lā} open? (The \emph{shahada} — \emph{lā ilāha illā llāh}.) Besides \textquotedblleft{}no,\textquotedblright{} what job does \emph{lā} do? (It \textbf{negates verbs} — \emph{lā aʿrif}, \textquotedblleft{}I don't know.\textquotedblright{})
\end{tcolorbox}
